\section{Results}
\label{sec:results}

\subsection{Accounting and Fresh-Answer Reliability}

The locked run accounts for all \TotalInitial{} planned initial cases, with
\TotalRepairs{} bounded repairs and \TotalRequests{} total request attempts.
There are \TotalFailed{} failed final cases and \TotalInitialErrors{} initial
request/format failures, including \TotalTruncated{} truncated outputs. The
journals record \ClientTimeouts{} client timeout attempts across initial and
repair phases. The matched experiment takes \MainHours{} hours on the
recorded hardware. Failed, empty, and partially answered cases remain in their
fixed denominators. Table~\ref{tab:reliability} reports case error together with
claim error, required-fact recall, and useful coverage for each source and method.
Each source/method cell contains 240 query checkpoints. The three failed requests
are transport ReadErrors in B0, one per source; they were retained without retries.

\begin{table}[t]
\caption{Fresh displayed answers by source. All rates are percentages. Recall
counts answerable required slots; coverage requires all such slots. An empty
answer contributes no checkable claims and reduces coverage. Errors are violations
of the bounded query contract, including interval scope.}
\label{tab:reliability}
\centering\small\setlength{\tabcolsep}{3.2pt}
\begin{tabular}{llrrrrr}
\toprule
Source & Method & Case err.\% & Claim err.\% & Recall\% & Coverage\% & Claims \\
\midrule
Synthetic & B0 & 52.9 & 43.2 & 78.8 & 60.0 & 588 \\
Synthetic & B1 & 0.0 & 0.0 & 79.3 & 60.5 & 335 \\
Synthetic & B2 & 0.0 & 0.0 & 79.3 & 60.5 & 336 \\
Synthetic & M1 & 0.0 & 0.0 & 79.3 & 60.5 & 336 \\
Synthetic & B3 & 0.0 & 0.0 & 79.3 & 60.5 & 335 \\
WESAD & B0 & 59.2 & 39.2 & 76.0 & 54.1 & 710 \\
WESAD & B1 & 0.0 & 0.0 & 76.9 & 55.9 & 440 \\
WESAD & B2 & 0.0 & 0.0 & 76.9 & 55.9 & 440 \\
WESAD & M1 & 0.0 & 0.0 & 76.9 & 55.9 & 440 \\
WESAD & B3 & 0.0 & 0.0 & 76.7 & 55.5 & 439 \\
PPG-DaLiA & B0 & 56.2 & 38.5 & 76.9 & 55.9 & 693 \\
PPG-DaLiA & B1 & 0.0 & 0.0 & 76.9 & 55.9 & 425 \\
PPG-DaLiA & B2 & 0.0 & 0.0 & 77.1 & 56.4 & 427 \\
PPG-DaLiA & M1 & 0.0 & 0.0 & 76.7 & 55.5 & 425 \\
PPG-DaLiA & B3 & 0.0 & 0.0 & 76.9 & 55.9 & 426 \\
\bottomrule
\end{tabular}

\end{table}

All checked methods have zero observed case and claim violations under the exact contract. B0 case violation rates are 52.9--59.2\%. M1 retains 76.7--79.3\% of answerable required facts, but supplies all answerable slots for only 55.5--60.5\% of eligible queries. The checked controls have closely similar recall and coverage.


The distinction between contract compliance and factual correctness materially
affects these results. A post-hoc descriptive check finds that
\OffTargetSupportedClaims{} of B0's \BaselineInvalidClaims{} invalid claims are
direct numerical background statements that pass the same exact checks at their
own stated interval. For example, a correctly labelled preceding-window EDA
median violates the target-interval rule. A comparison declaring that statement
as a parent can consequently fail the dependency check despite correct arithmetic.
Table~\ref{tab:contractscope} gives source-wise counts. The frozen primary scores
are unchanged; their error reductions must not be read as equivalent reductions
in hallucinations, and their recall also depends on this contract.

These comparisons describe the final displayed content after the common bounded
repair policy. A low error rate cannot alone establish useful answering: removing
an unsupported sentence can lower error while leaving a required value absent.
Figure~\ref{fig:reliability} places error and recall on the same source-wise
comparison. Table~\ref{tab:failures} exposes abstention, reference validity,
repairs, and failures. Error-component counts and answerability-stratified partial
answer/abstention counts are available in the reproducible analysis tables.

\begin{figure}[t]
\centering
\includegraphics[width=\textwidth]{minimum_v1_reliability.pdf}
\caption{Query-contract case violations (top) and required-fact recall (bottom) for the three
waveform sources. The five bars within each panel use the same evidence and
generation settings. Lower error is desirable only alongside retained recall.
These are aggregate proportions; paired subject-cluster intervals are reported
for the prospective contrasts in Table~\ref{tab:contrasts}.}
\label{fig:reliability}
\end{figure}

\subsection{Required Changes and Retained Content}

The exact symbolic fixtures isolate maintenance from generation. They contain
two claims requiring change after a correction or invalidation, an independently
supported OR claim, and an unaffected baseline claim. Benign and irrelevant
updates test stability. Table~\ref{tab:fixture} aggregates these four constructed
interventions. Full graph and full relational maintenance identify both direct
and downstream required changes and retain alternative support. Direct-only
maintenance identifies the direct change while leaving the downstream old claim
active. Restoration is also checked by replaying a third version. These are
deterministic correctness cases, not independent participant observations, and
no sampling interval is attached to them.

\begin{table}[t]
\caption{Common symbolic fixture. Required changes occur in correction and
invalidation modes; benign and irrelevant modes contribute unaffected claims.
The identical candidate structure is replayed into every method.}
\label{tab:fixture}
\centering\small
\begin{tabular}{lrrr}
\toprule
Method & Corrected/required & Residual & Collateral/unaffected \\
\midrule
B0 & 0/4 & 4 & 0/12 \\
B1 & 2/4 & 2 & 0/12 \\
B2 & 2/4 & 2 & 0/12 \\
M1 & 4/4 & 0 & 0/12 \\
B3 & 4/4 & 0 & 0/12 \\
\bottomrule
\end{tabular}

\end{table}

Generated persistent displays have a different denominator: the method must
first produce a valid claim that later requires change. Table~\ref{tab:correction}
therefore reports the actual corrected/required and collateral/unaffected counts,
as well as residual incorrect active content. Figure~\ref{fig:correction} displays
the corresponding proportions. These are claim-checkpoint opportunities, and
their number can differ across methods because initial answers differ. Initially
incorrect content is tracked separately rather than credited as correctable
content that the method once supported validly.

All four checked methods correct every eligible generated claim and have no
collateral withdrawals in this run. Thus, the waveform task supplies no measured
correction advantage for transitive propagation over direct checking. B0 retains
every eligible outdated claim. The symbolic and generated results answer different
questions and should not be pooled.

\begin{table}[t]
\caption{Persistent generated displays. Completeness is conditional on an initially
valid claim requiring change. Residual counts those still active; collateral
counts withdrawal of content still valid. Undefined rates are shown as dashes.}
\label{tab:correction}
\centering\small\setlength{\tabcolsep}{3pt}
\begin{tabular}{llrrrr}
\toprule
Source & Method & \shortstack{Corrected/\\required} & \shortstack{Complete\\(\%)} & Residual & \shortstack{Collateral/\\unaffected} \\
\midrule
Synthetic & B0 & 0/71 & 0.0 & 71 & 0/245 \\
Synthetic & B1 & 73/73 & 100.0 & 0 & 0/247 \\
Synthetic & B2 & 73/73 & 100.0 & 0 & 0/247 \\
Synthetic & M1 & 73/73 & 100.0 & 0 & 0/247 \\
Synthetic & B3 & 73/73 & 100.0 & 0 & 0/247 \\
WESAD & B0 & 0/71 & 0.0 & 71 & 0/336 \\
WESAD & B1 & 74/74 & 100.0 & 0 & 0/349 \\
WESAD & B2 & 74/74 & 100.0 & 0 & 0/349 \\
WESAD & M1 & 74/74 & 100.0 & 0 & 0/349 \\
WESAD & B3 & 73/73 & 100.0 & 0 & 0/348 \\
PPG-DaLiA & B0 & 0/80 & 0.0 & 80 & 0/326 \\
PPG-DaLiA & B1 & 80/80 & 100.0 & 0 & 0/329 \\
PPG-DaLiA & B2 & 80/80 & 100.0 & 0 & 0/330 \\
PPG-DaLiA & M1 & 80/80 & 100.0 & 0 & 0/328 \\
PPG-DaLiA & B3 & 80/80 & 100.0 & 0 & 0/330 \\
\bottomrule
\end{tabular}

\end{table}

\begin{figure}[t]
\centering
\includegraphics[width=\textwidth]{minimum_v1_correction.pdf}
\caption{Persistent-display correction completeness (top) and collateral withdrawal
(bottom). Denominators appear in Table~\ref{tab:correction}. A zero collateral
proportion has a different interpretation from an undefined correction rate;
undefined bars are labeled n/a. These are replay outcomes, not open-loop
arrival-to-replacement measurements.}
\label{fig:correction}
\end{figure}

M1 emits \GraphTotalClaims{} displayed claims, of which \GraphParentClaims{}
declare parent claims. The bounded numerical contract also requires direct
evidence references. This structure constrains the opportunity for transitive
maintenance to add value over direct checking in the waveform task. The common
fixture demonstrates a deliberately necessary transitive path; its result should
not be substituted for an empirical advantage across unrestricted generated
reasoning chains. No generated persistent claim has a declared parent dependency
as its sole later exact-check error in this run.

\subsection{Paired Contrasts}

Table~\ref{tab:contrasts} reports the two prospective principal comparisons by
source. Differences are M1 minus the control in percentage points, computed
within subjects. Negative case-error differences favor M1; positive recall and
correction differences favor M1. The table keeps these outcomes separate so
error reductions cannot hide lost information. Conditional correction contrasts
disclose their eligible subject and episode counts. Pointwise bootstrap intervals quantify
participant-level variability within this fixed generation run; they do not
measure variability across model seeds or remove shared measurement error.
The zero differences and zero-width intervals for the principal correction and
case-error contrasts arise because all observed paired subject differences are
identical. They do not establish universal equivalence or a zero error probability.

\begin{table}[t]
\caption{Principal paired differences with 2,000 subject-cluster resamples.
The M1--B1 reliability contrast is interpreted jointly for error and recall.
N (S/E) gives retained subjects/episodes.
Intervals are pointwise 95\%; no hypothesis-test $p$-values are reported.}
\label{tab:contrasts}
\centering\small\setlength{\tabcolsep}{3pt}
\begin{tabular}{lllrrl}
\toprule
Source & Contrast & Outcome & N (S/E) & Difference (pp) & 95\% interval (pp) \\
\midrule
Synthetic & M1--B2 & Correction & 10/20 & 0.00 & [0.00, 0.00] \\
Synthetic & M1--B1 & Case error & 10/20 & 0.00 & [0.00, 0.00] \\
Synthetic & M1--B1 & Fact recall & 10/20 & 0.00 & [0.00, 0.00] \\
WESAD & M1--B2 & Correction & 10/20 & 0.00 & [0.00, 0.00] \\
WESAD & M1--B1 & Case error & 10/20 & 0.00 & [0.00, 0.00] \\
WESAD & M1--B1 & Fact recall & 10/20 & 0.00 & [0.00, 0.00] \\
PPG-DaLiA & M1--B2 & Correction & 10/20 & 0.00 & [0.00, 0.00] \\
PPG-DaLiA & M1--B1 & Case error & 10/20 & 0.00 & [0.00, 0.00] \\
PPG-DaLiA & M1--B1 & Fact recall & 10/20 & -0.25 & [-0.75, 0.00] \\
\bottomrule
\end{tabular}

\end{table}

All checked conditions use the same fresh-answer checks and repair policy on
matched evidence. Small M1--B1 differences therefore do not isolate an additional
fresh-answer mechanism; persistent displays supply the maintenance comparison.
The M1--B3 differences in Table~\ref{tab:relational} compare generated outcomes
under equivalent support semantics. Fresh outputs can still differ under the
same requested seed because scheduling and batching are not bitwise identical.
The deterministic fixtures and systems assessment hashes supply the separate
test of logical equivalence. Differences in generated answers must not be
attributed to graph storage alone.

\subsection{Open-Loop Latency and Resource Use}

All six systems cells complete, and graph/relational logical assessment hashes
agree for every offered rate. Table~\ref{tab:streaming} reports flag latency,
time until candidate completion, deadline misses, and peak queues from arrival.
Figure~\ref{fig:streaming} displays the load dependence. The M1 implementation
misses \GraphFlagMisses{} of \SystemsRepairTargets{} flag targets and
\GraphRepairMisses{} of \SystemsRepairTargets{} adequate-replacement targets.
Replacement failure includes inadequate content as well as lateness. Fast
withdrawal therefore does not imply a useful replacement within five seconds.
All \SystemsExplanations{} systems displays contain a single target observation
and omit the required difference, even after the bounded repair attempt.
Successful-replacement latency is therefore undefined; the reported candidate
latency measures when these incomplete attempts finish.

At one and five events/s, both backends meet every flag deadline, and the
relational implementation has lower p95 flag latency. At 20 events/s both
accumulate event and answer queues and miss flag deadlines. M1 has lower p95
flag and candidate latency at that rate, but more flag misses. This reversal
across summaries cautions against a single storage-engine ranking.

\begin{table}[t]
\caption{Open-loop systems results. Each cell offers events for 60 seconds and
then drains. Q columns are peak queued events and pending explanations. Flag
and repair columns show misses/opportunities for one- and five-second targets.
Candidate latency includes inadequate outputs; no complete replacement is observed.}
\label{tab:streaming}
\centering\small\setlength{\tabcolsep}{2.4pt}
\begin{tabular}{rlrrrrrr}
\toprule
Events/s & Method & \shortstack{Flag p95\\(ms)} & \shortstack{Candidate\\p95 (s)} & \shortstack{Flag\\misses} & \shortstack{Repair\\misses} & \shortstack{Event\\Q} & \shortstack{Answer\\Q} \\
\midrule
1 & M1 & 43.9 & 11.0 & 0/2 & 2/2 & 0 & 2 \\
1 & B3 & 3.7 & 10.9 & 0/2 & 2/2 & 0 & 2 \\
5 & M1 & 52.8 & 18.7 & 0/8 & 8/8 & 0 & 7 \\
5 & B3 & 23.7 & 19.5 & 0/8 & 8/8 & 0 & 7 \\
20 & M1 & 33,687.0 & 218.7 & 21/30 & 30/30 & 347 & 62 \\
20 & B3 & 77,698.4 & 237.9 & 17/30 & 30/30 & 394 & 50 \\
\bottomrule
\end{tabular}

\end{table}

\begin{figure}[t]
\centering
\includegraphics[width=\textwidth]{streaming.pdf}
\caption{Offered load versus p95 flag and candidate-completion latency and peak
queues for M1 and B3. Dotted lines mark the one-second flag and five-second
replacement targets. The flag axis is logarithmic. The independent producer
continues scheduling events while requests wait or generate; measurements
include the resulting backlog.}
\label{fig:streaming}
\end{figure}

The frozen adequacy check uses the request snapshot. A descriptive check also
remaps evidence to the latest ingested versions at completion
(Table~\ref{tab:staleness}), preserving the original scores. With no adequate
input-snapshot answers, zero answers ``became stale''; this is uninformative about
preserving useful replacements. Rechecking before publication and coalescing
obsolete requests remain unevaluated scheduling policies.

The experiment uses one RTX 5090; Table~\ref{tab:hardware} gives hardware and model
details. The development workload with one concurrent request has median and p95 generation
latencies of \PilotInteractivePFifty{} and \PilotInteractivePNinetyFive{} seconds, including
bounded repairs. The separate batched pilot completes \PilotBatchRate{} initial
cases per second, including repair and database cost. These are different workloads;
batch throughput does not establish low interactive latency.
The primary worker verifies all parameters and first-forward tensors on CUDA,
with no CPU weight offload. The systems track makes \SystemsCalls{} requests,
within its separate 472-request ceiling. Table~\ref{tab:cost} includes retained
development failures, final experiments, and the verifier. Stage wall times
include host work and idle gaps and are not CUDA kernel times. One-second
telemetry measures observed VRAM and utilization; cgroup memory includes file
cache and other processes in the pod. Database node, edge, chain-length and
physical-size accounting is included in the artifact report. Physical Neo4j
storage also contains development scopes, so its size is not presented as an
equal-content compression comparison. No hourly account price was available;
we report time and tokens without an invented monetary cost.

\begin{table}[t]
\caption{Measured request and token accounting. Pilot rows include their separate
interactive workload. Systems pilot v1 includes requests from its failed scoring
smoke test. ``No usage'' counts attempts without returned token counts: their
server-side cost is unknown, not zero. Token sums include available usage only.
Wall time excludes package/model downloads and manuscript work.}
\label{tab:cost}
\centering\small\setlength{\tabcolsep}{2.8pt}
\begin{tabular}{lrrrrr}
\toprule
Workload & Requests & No usage & Input tokens & Output tokens & Wall min \\
\midrule
Matched test & 5,220 & 3 & 8,525,358 & 2,880,024 & 173.2 \\
pilot v1 & 207 & 0 & 256,281 & 103,092 & 11.4 \\
pilot v2 & 164 & 0 & 230,554 & 90,664 & 9.6 \\
pilot v3 & 186 & 0 & 300,657 & 106,302 & 9.3 \\
pilot v4 & 187 & 0 & 302,546 & 106,911 & 9.3 \\
Systems test & 472 & 0 & 416,220 & 206,856 & 15.1 \\
Automated audit & 64 & 0 & 143,213 & 2,786 & 0.6 \\
Systems pilot v1 & 8 & 0 & 6,792 & 3,060 & 0.4 \\
Systems pilot v2 & 8 & 0 & 6,792 & 3,060 & 0.4 \\
\bottomrule
\end{tabular}

\end{table}

\subsection{Automated Fidelity Audit and Review Demonstration}

The separate verifier parses \AuditParsed{} of 60 sampled nonempty final
explanations, flags \AuditFlagged{}, and disagrees with the exact claim checks
on \AuditDisagreements{}. It gives the expected result on \AuditCalibration{}
of four diagnostic examples. These counts describe an auxiliary automated audit,
not a validated estimate of semantic accuracy. Selection is stratified across
methods and sources and conditional on a nonempty checkpoint-2 display. Empty
outputs are already represented in coverage and abstention metrics. Raw audit
prompts and responses permit inspection of both flags and apparent false alarms.
The verifier accepts \AuditAcceptedInvalid{} exact-invalid outputs, including
incorrect differences and absent citations as well as interval-contract violations.
It flags \AuditFlaggedValid{} exact-valid outputs; its rationales include treating
a null measurement as present and rejecting rounding within the stated tolerance.
These visible failures limit its value as a semantic check despite passing the
four simple diagnostics.

The local review dashboard demonstrates the maintenance operations in an isolated
workspace. An issue queue prioritizes explicit evidence problems and downstream
claim count; a selected claim exposes its local graph, source feature values,
time windows, and versions. Reviewers can reject evidence, accept a proposed
numerical correction, or mark an interpretation unresolved. Actions append
ReviewEvents and explanation versions and show before/after text. Scripted
examples demonstrate dependent-statement withdrawal, retained OR support, and
a mistaken rejection that removes valid content. These examples establish
functionality and an observable limitation of review actions; no human performance
or review-speed outcome was measured.
